\documentclass[11pt,a4paper]{article}

% ── FONTS ────────────────────────────────────────────────────
\usepackage{fontspec}
\setmainfont{Calibri}
\setsansfont{Calibri}

% ── LAYOUT ───────────────────────────────────────────────────
\usepackage[top=1.8cm,bottom=2cm,left=2cm,right=2cm]{geometry}
\usepackage{graphicx}
\usepackage{xcolor}
\usepackage{booktabs}
\usepackage{tabularx}
\usepackage{enumitem}
\usepackage{setspace}
\usepackage{fancyhdr}
\usepackage{titlesec}
\usepackage{hyperref}
\usepackage{microtype}
\usepackage{array}
\usepackage{colortbl}

\hypersetup{colorlinks=true,linkcolor=navyblue,urlcolor=accentteal,citecolor=navyblue}

% ── BRAND COLORS ────────────────────────────────────────────
\definecolor{navyblue}{HTML}{003566}
\definecolor{gold}{HTML}{C9A84C}
\definecolor{lightgray}{HTML}{F4F6FB}
\definecolor{darkgray}{HTML}{4A4A4A}
\definecolor{accentteal}{HTML}{0077B6}

% ── SECTION STYLING ─────────────────────────────────────────
\titleformat{\section}[block]
  {\normalfont\bfseries\fontsize{11}{13}\selectfont\color{navyblue}}
  {}{0em}{}
\titlespacing*{\section}{0pt}{10pt}{4pt}

\titleformat{\subsection}[block]
  {\normalfont\bfseries\fontsize{10}{12}\selectfont\color{accentteal}}
  {}{0em}{}
\titlespacing*{\subsection}{0pt}{6pt}{2pt}

% ── HEADER / FOOTER ─────────────────────────────────────────
\setlength{\headheight}{16pt}
\pagestyle{fancy}
\fancyhf{}
\fancyhead[L]{\footnotesize\color{navyblue}\textbf{AIGGPA Bhopal} \quad\textbar\quad
  Human Capital: Investment in Education, Skill Development and Employment}
\fancyhead[R]{\footnotesize\color{navyblue}April 2026}
\fancyfoot[C]{\footnotesize\color{navyblue}\thepage}
\renewcommand{\headrulewidth}{0.4pt}
\renewcommand{\headrule}{\color{navyblue}\hrule width\headwidth height\headrulewidth}
\renewcommand{\footrulewidth}{1pt}
\renewcommand{\footrule}{\color{gold}\hrule width\headwidth height 1pt}

% ── LISTS ───────────────────────────────────────────────────
\setlist[itemize]{leftmargin=1.2em, itemsep=1pt, topsep=2pt}

\newcommand{\rupee}{\textrm{\fontspec{Calibri}₹}}

% ────────────────────────────────────────────────────────────
\begin{document}
\onehalfspacing

\begin{center}
  {\fontsize{12}{14}\selectfont\bfseries\color{navyblue}
  Skilling vs. Degree: An Analytical Overview of\\[2pt]
  Demand, Gaps \& the Changing Technology Landscape}\\[6pt]
  {\small\color{darkgray} Himani \quad\textbar\quad AIGGPA Summer Intern 2026}
\end{center}

\vspace{0.2cm}

\section{1. The Degree--Skill Paradox in India}

India produces over \textbf{10 million graduates annually}, yet the \textit{India Skills Report 2025} (CII--Wheebox) reveals that only \textbf{54.81\%} are considered employable by industry standards -- up marginally from 51.25\% in 2024. This ``degree-rich, skill-poor'' paradox arises because:

\begin{itemize}
  \item Traditional curricula emphasise theoretical knowledge over applied competency.
  \item Over \textbf{50\% of graduates} work in roles below their qualification level (underemployment).
  \item Employers increasingly seek demonstrable skills in \textbf{AI/ML, Cloud Computing, Cybersecurity, and Data Analytics} rather than degree credentials alone.
  \item The hiring paradigm is shifting: companies like Google, IBM, and Infosys now offer ``degree-optional'' roles, prioritising certifications and portfolios.
\end{itemize}

\vspace{0.15cm}
\section{2. Current Demand \& Skill Gaps -- National Picture}

\renewcommand{\arraystretch}{1.25}
\begin{table}[h!]
\centering
\footnotesize
\caption{\textbf{Employability \& Demand Snapshot (India Skills Report 2025)}}
\vspace{0.1cm}
\begin{tabularx}{\textwidth}{>{\columncolor{lightgray}}l >{\centering\arraybackslash}X >{\centering\arraybackslash}X}
\toprule
\rowcolor{navyblue}
\textcolor{white}{\textbf{Category}} & \textcolor{white}{\textbf{Employability Rate}} & \textcolor{white}{\textbf{Key Gap Area}} \\
\midrule
Management Graduates (MBA)   & 78\%    & Soft skills, leadership \\
\rowcolor{white}
Engineering Graduates (B.Tech) & 71.5\% & AI/ML, hands-on coding \\
MCA Graduates                 & 71\%    & Cloud, DevOps \\
\rowcolor{white}
Science Graduates (B.Sc.)      & 58\%    & Industry exposure, analytics \\
Arts/Commerce Graduates        & $<$45\% & Digital literacy, communication \\
\bottomrule
\end{tabularx}
\vspace{0.05cm}
\par\scriptsize Source: India Skills Report 2025, CII--Wheebox--AICTE; NASSCOM Academia Survey 2025.
\end{table}

\textbf{High-demand technology areas (2025--26):} Artificial Intelligence, Generative AI, Machine Learning, Cloud Architecture (AWS/Azure/GCP), Cybersecurity, Full-Stack Development, Data Engineering, and Semiconductor Design.

\textbf{Critical soft-skill gaps:} Communication, critical thinking, teamwork, and adaptability -- skills that neither degrees nor technical certifications alone can guarantee.

\vspace{0.15cm}
\section{3. Madhya Pradesh -- State-Level Gaps \& Initiatives}

MP faces a sharper version of the national challenge:

\begin{itemize}
  \item \textbf{Low vocational training:} Only $\sim$\textbf{1.2\%} of MP's 15+ population has received formal vocational training, compared to the national average of 2.4\% (NSSO/PLFS data).
  \item \textbf{29 lakh+ registered job seekers} on the state employment portal (March 2025), many holding degrees but lacking industry-ready skills.
  \item \textbf{Emerging demand:} MP's new IT/ITeS Policy, ESDM Policy, and GCC Policy 2025 are attracting tech investment to Indore, Bhopal, and Gwalior -- creating demand for skills the current education system does not adequately supply.
\end{itemize}

\textbf{State Responses:}
\begin{itemize}
  \item \textit{Mukhyamantri Seekho-Kamao Yojana} -- ``Learn and Earn'' scheme combining on-the-job training with stipends.
  \item \textit{Global Skills Park (GSP), Bhopal} -- ADB-funded centre of excellence providing internationally benchmarked, advanced vocational training in emerging trades.
  \item \textit{ITI Modernisation} -- MP achieved a Top-5 national ranking in ITI performance (2025); curricula being upgraded to include AI, IoT, and electric vehicle technology modules.
  \item Industry partnerships with Reliance, Trident, and Jindal for placement-linked training.
\end{itemize}

\vspace{0.15cm}
\section{4. Key Takeaway}

The evidence is clear: \textbf{degrees open doors, but skills keep them open}. In a technology landscape disrupted by AI and automation, India -- and MP in particular -- must accelerate the shift from a ``degree-first'' to a ``skill-first'' human capital model. The Union Budget 2026--27's 62\% increase in skill development allocation (\rupee9,886 Cr) and MP's own skilling infrastructure investments signal that this transition is underway, but the pace must match the urgency of a 370-million-strong youth population entering the workforce.

\vspace{0.25cm}
\noindent\rule{\textwidth}{0.5pt}
{\scriptsize\textbf{References:}
(1)~India Skills Report 2025, CII--Wheebox--AICTE, wheebox.com.
(2)~NASSCOM, \textit{Academia--Industry Alignment Survey 2025}, nasscom.in.
(3)~Ministry of Skill Development \& Entrepreneurship, \textit{Annual Report 2024--25}, msde.gov.in.
(4)~Asian Development Bank, \textit{MP Skills Development Project}, adb.org.
(5)~Government of Madhya Pradesh, \textit{GCC Policy 2025 \& IT/ITeS Policy}, mpinfo.org.
(6)~NSSO/PLFS, \textit{Periodic Labour Force Survey 2023--24}, mospi.gov.in.
}

\end{document}
